% this file is called up by thesis.tex
% content in this file will be fed into the main document

%: ----------------------- name of chapter  -------------------------

\chapter{Literature Review} % top level followed by section, subsection

%: ----------------------- paths to graphics ------------------------

% change according to folder and file names
\ifpdf
    \graphicspath{{2_chapter2/figures/PNG/}{2_chapter2/figures/PDF/}{2_chapter2/figures/}}
\else
    \graphicspath{{2_chapter2/figures/EPS/}{2_chapter2/figures/}}
\fi

%: ----------------------- contents from here ------------------------

\section{Introduction} 
\begin{itemize}
    \item Purpose of this chapter
    \item Brief overview on TILDA as a unique dataset (rich, longitudinal, multi-domai...)
    \item Scope: frailty, mental health, chronic disease, multimorbidity, nutrition, and ML models
\end{itemize}

\section{Chronic Disease and Ageing}
\begin{itemize}
    \item The burden of chronic conditions in older adults
    \item Healthy ageing framework \gls{WHO}
    \item Importance of multimorbidity, frailty, mental health
\end{itemize}

\subsection{Frailty and Physical Function}
\begin{itemize}
    \item Gait speed, grip strength, functional mobility decline 
    \item Predictors of frailty transitions (\cite{mccarthyAssociationMetabolicSyndrome2023}, \cite{huangAgeModifyAssociations2019}, \cite{oconnellOrthostaticHypotensionOrthostatic2015})
\end{itemize}

\subsection{Mental Health in Older Adults}
\begin{itemize}
    \item Depression, anxiety, loneliness, social connectedness (\cite{santiniProtectivePropertiesActBelongCommit2017}, \cite{Symposia2011}, \cite{briggsWishLaterLife2021})
    \item Links between physical activity and depression (\cite{mcdowellAssociationsPhysicalActivity2018}, \cite{lairdPhysicalActivityDose2023}, \cite{lairdPhysicalActivityDepression2023})
\end{itemize}

\subsection{Nutrition and Biological Ageing}
\begin{itemize}
    \item Vitamin D, lutein, zeaxanthin, \gls{MPOD} (\cite{murphyPlasmaLuteinZeaxanthin2023}, \cite{mccarthyAssociationVitaminDeficiency2022}, lairdVitaminDeficiencyIreland2020, feeneyLowMacularPigment2013)
    \item Malnutrition and its predictors (\cite{bardonPredictorsIncidentMalnutrition2020})
    \item Biological vs chronological ageing (\cite{mccarthyMetabolicSyndromeAccelerates2023}, \cite{aziziEvidenceAssociationAllostatic2024})
\end{itemize}

\section{Machine Learning in Ageing and Chronic Disease}
\begin{itemize}
    \item Overview: Why ML is valuable in ageing research \cite{zhangMachineLearningModels2025}
    \item Transition from regression-based models → ensemble ML → deep learning(cite a example here)
\end{itemize}

\subsection{Feature Sets in ML Studies}
\begin{itemize}
    \item Clinical, sociodemographic, lifestyle, cognition, psychosocial
    \item Over-reliance on physical health predictors vs underuse of nutrition/social variables
\end{itemize}

\subsection{ML Models Used}
\begin{itemize}
    \item Logistic regression, Markov models, \gls{LDA} (traditional)
    \item Ensembles \gls{RF}, \gls{GBM}, \gls{DNN} – e.g., (\cite{xuTailoredMachineLearning2023}) diabetes prediction
    \item Unsupervised clustering for dementia (\cite{bayeneleonoreUnsupervisedMachineLearning2018})
    \item Brain age prediction models (\cite{boyleBrainpredictedAgeDifference2021}, \cite{a.shirsathSlowerSpeedBlood2024})
\end{itemize}

\subsection{Performance and Validation}
\begin{itemize}
    \item \gls{AUC}, R² ranges across tasks (diabetes, gait speed, mortality, depression)
    \item Limitations: internal validation only, cross-sectional data, no generalisation
\end{itemize}

\subsection{Explainability in ML for Ageing}
\begin{itemize}
    \item \gls{SHAP}, TreeExplainer examples (\cite{xuTailoredMachineLearning2023}, \cite{davisLinearRegressionbasedMachine2021}, \cite{davisComparisonGaitSpeed2021})
    \item Challenges of black-box models in clinical practice
\end{itemize}

\section{Gaps in the Literature}
\begin{itemize}
    \item Limited population diversity (mostly Irish, 50+)
    \item Over-reliance on physical predictors; neglect of nutrition, social, religious/spiritual, environmental variables
    \item Dominance of regression models; limited advanced ML such as neural networks, transformers
    \item Validation gaps – very few external validations, especially for ensemble and brain-age models
    \item Underexplored outcomes: pain, sleep, QoL
    \item Cross-sectional bias; weak time-series modelling
    \item Poor transparency on missing data handling
    \item Limited integration of nutrition biomarkers
    \item Lack of biological age prediction vs chronological
    \item Interpretability challenges in ML
\end{itemize}

\section{Chapter Summary}
\begin{itemize}
    \item I need to recap the main findings from LR systematic mapping excel file
    \item Identify your contribution niche:
    \begin{itemize}
        \item integrating nutrition + lifestyle + clinical predictors
        \item ML + explainability
        \item time-based prediction (longitudinal TILDA data)
    \end{itemize}
    \item Lead into Methodology (Chapter 3)
\end{itemize}



% ---------------------------------------------------------------------------
%: ----------------------- end of thesis sub-document ------------------------
% ---------------------------------------------------------------------------

